\chapter{RESEARCH METHODOLOGY}

This chapter details the research methodology employed to develop the Multiple Embedding Steganography (MES) framework. It outlines the research design, the proposed parallel embedding algorithm, and the instruments used for data collection and analysis.

% =========================================================
\section{Research Design}
% =========================================================
This research follows an experimental development research design. The core of the study is the design and implementation of a new steganographic algorithm that utilizes parallel processing in the RGB color space. The development process involves:
\begin{enumerate}
    \item \textbf{Algorithm Design:} Conceptualizing the Parallel Embedding architecture.
    \item \textbf{Implementation:} Coding the algorithm using a high-level programming language.
    \item \textbf{Experimentation:} Testing the algorithm on standard dataset images.
    \item \textbf{Evaluation:} Analyzing performance using quantitative metrics (PSNR, bpp).
\end{enumerate}

% =========================================================
\section{Proposed Method: Parallel Embedding Architecture}
% =========================================================
The proposed MES framework diverges from traditional sequential methods by treating the Red, Green, and Blue channels as independent data carriers.

\subsection{Channel Decomposition}
The first step involves decomposing the RGB cover image $I$ into its three constituent channels: $I_R$, $I_G$, and $I_B$.
\begin{equation}
    I(x,y) \rightarrow \{ I_R(x,y), I_G(x,y), I_B(x,y) \}
\end{equation}

\subsection{Edge Detection and Classification}
For each channel, an edge detection mechanism (using the Canny operator) is applied to classify pixels into "Edge" or "Smooth" regions. 
\begin{itemize}
    \item \textbf{Edge Regions:} High variance. Higher embedding capacity is assigned here as the human eye is less sensitive to modifications.
    \item \textbf{Smooth Regions:} Low variance. Lower embedding capacity is used to preserve visual quality.
\end{itemize}

\subsection{Parallel Embedding Process}
The secret data $S$ is split into three segments $S_R, S_G, S_B$ proportional to the capacity of each channel. These segments are embedded simultaneously:
\begin{itemize}
    \item $S_R \rightarrow I_R$ using Adaptive PVD
    \item $S_G \rightarrow I_G$ using Adaptive PVD
    \item $S_B \rightarrow I_B$ using Adaptive PVD
\end{itemize}
This parallelization ensures that embedding in the Blue channel is not dependent on the Red channel, preventing error propagation.

% =========================================================
\section{Experiment Scenario}
% =========================================================
The experiments are conducted using a standard set of test images widely used in image processing research, including "Lenna", "Baboon", "Peppers", and "Airplane". Each image is $512 \times 512$ pixels.
The secret data used for embedding includes both random text data and full-size $512 \times 512$ RGB images to test the high-capacity claim.

% =========================================================
\section{Tools for Data Analysis}
% =========================================================
To evaluate the performance of the proposed method, the following metrics are used:

\subsection{Peak Signal-to-Noise Ratio (PSNR)}
PSNR measures the quality of the stego-image compared to the original cover image. A higher PSNR indicates better quality.
\begin{equation}
    PSNR = 10 \cdot \log_{10} \left( \frac{MAX_I^2}{MSE} \right)
\end{equation}

\subsection{Embedding Capacity}
Capacity is measured in bits per pixel (bpp). The target is a capacity $> 3.0$ bpp, which theoretically allows embedding a full color image within another.
